\documentclass[11pt]{article}

% ---- Preamble: safe math & fonts ----
\usepackage[T1]{fontenc}
\usepackage[utf8]{inputenc}
\usepackage{lmodern}
\usepackage{geometry}
\geometry{margin=1in}
\usepackage{amsmath, amssymb, amsbsy, amsthm, mathtools}
\usepackage{graphicx}
\usepackage{enumitem}
\usepackage{hyperref}
\hypersetup{colorlinks=true,linkcolor=blue,urlcolor=blue,citecolor=blue}

% Convenience macros
\newcommand{\E}{\mathbb{E}}
\newcommand{\KL}{\mathrm{KL}}
\newcommand{\given}{\,\mid\,}
\newcommand{\R}{\mathbb{R}}

\title{Study Guide --- \emph{From Behaviorism to Language Models:}\\
A Formal Correspondence Between the Matching Law and Token Prediction}
\author{Companion to the preprint by Wesley R. Elsberry}
\date{\today}

\begin{document}
\maketitle

\section*{Purpose and Audience}
\textbf{Audience:} graduate students and advanced undergraduates in psychology, cognitive science, machine learning, and adjacent fields.\\
\textbf{Goal:} make the preprint maximally accessible by explaining background, acronyms, key equations, theorem derivations, and the broader research agenda.

\section{How to Use This Guide}
\begin{enumerate}[leftmargin=*]
  \item Read the Background Primer to level-set on behaviorism, LLM training, and the core math objects (cross-entropy, KL, logits).
  \item Walk through the Theorem Roadmap step-by-step; each theorem includes assumptions, a derivation sketch, intuition, and pitfalls.
emacs merge mode  \item Consult the Glossary whenever acronyms or terms appear.
  \item Use the Micro-Exercises to build intuition with simple numbers.
  \item Review the Broader Research Program (Table~\ref{tab:program}) for project ideas and empirical tests.
  \item See Appendix~\ref{app:ngram} in the preprint for code and examples contrasting $n$-gram “parrots” with LLMs.
\end{enumerate}

\section{Background Primer}

\subsection{Behaviorism \& the Matching Law (Operant Conditioning)}
\textbf{Operant conditioning} holds that behavior is shaped by its consequences (reinforcement schedules). \textbf{The matching law} (Herrnstein) states that allocation of responses across options tends to match the allocation of reinforcement across those options: if option A yields $70\%$ of rewards and option B $30\%$, responses tend toward $70{:}30$ (\emph{strict matching}). \textbf{Generalized matching} (Baum) summarizes deviations via two parameters:
\begin{itemize}[leftmargin=*]
  \item \emph{Sensitivity} (slope): how strongly responses track reinforcement differences.
  \item \emph{Bias} (intercept): systematic preference independent of reinforcement.
\end{itemize}


\subsection{LLMs: Training, Alignment, Inference}
\textbf{Pretraining (Maximum Likelihood / Cross-Entropy):} learn to assign high probability to observed next tokens, i.e., minimize the divergence between model distribution $\pi_\theta(\cdot\given c)$ and corpus distribution $q(\cdot\given c)$.\\
\textbf{Alignment (RLHF / DPO):} introduce preference/reward signals and a KL penalty that keeps the updated policy close to a reference policy.\\
\textbf{Inference (Decoding):} use \emph{temperature} and related controls to modulate sampling variability.

\subsection*{Preview: Two Kinds of Grounding}

In the preprint we introduce two complementary terms to clarify how learning can
be said to be ``grounded.'' \emph{Allocation-grounding} refers to lawful,
substrate-independent allocation of responses in proportion to contingencies,
captured formally by the matching law. \emph{Gestalt-grounding} refers to the
particular experiential channel through which contingencies are encountered:
sensory-perceptual for organisms, textual-corpus for LLMs. Readers should note
that the full treatment appears in the Discussion, but keeping this distinction
in mind will help interpret the proofs and examples that follow.

\subsection*{What Does ``Lawful'' Mean in Behavior Analysis?}

Skinner often described behavior as ``lawful,'' but the term can be confusing
to modern readers. Here, ``lawful'' does not mean legal, moral, or normative. It
means that behavior can be described by regular, predictable relations to the
environment, much like how physics describes the lawful relation between force
and acceleration. 

For example, the matching law is a lawful relation: when the reinforcement rate
on one option is doubled, the proportion of responses allocated to that option
predictably increases. Describing such relations as lawful highlights that
behavior is not random or mysterious but can be analyzed, modeled, and even
predicted under the right conditions.



\section{Theorem Roadmap: Derivations, Intuition, Pitfalls}

\subsection*{Notation}
Contexts $c$, tokens $y$, corpus conditionals $q(\cdot\given c)$, model (policy) $\pi_\theta(\cdot\given c)$, reference policy $\pi_{\mathrm{ref}}$, reward $r(x,y)$ for alignment inputs $x$.

\subsection{Theorem 1: Strict Matching under MLE}
\textbf{Claim.} If $q(\cdot\given c)$ is representable by the model class, minimizing cross-entropy yields $\pi^*(\cdot\given c)=q(\cdot\given c)$. Thus for any two continuations $a,b$,
\begin{equation*}
\frac{\pi^*(a\given c)}{\pi^*(b\given c)}=\frac{q(a\given c)}{q(b\given c)}.
\end{equation*}
\noindent\emph{(strict matching)}

\paragraph{Assumptions.}
(1) Realizable model class; (2) global optimum is reached.

\paragraph{Derivation (3 steps).}
\begin{align*}
\mathcal{L}_{\mathrm{MLE}}
&= \E_{(c,y)\sim q}\big[-\log \pi_{\theta}(y\given c)\big]
= \E_{c\sim q}\Big[\,H\big(q(\cdot\given c),\,\pi_{\theta}(\cdot\given c)\big)\Big]\\
&= \E_{c\sim q}\Big[\,H\big(q(\cdot\given c)\big) + \KL\!\big(q(\cdot\given c)\,\Vert\,\pi_{\theta}(\cdot\given c)\big)\Big].
\end{align*}
$H(q)$ does not depend on $\theta$, so minimization forces $\KL(q\Vert \pi_\theta)\to 0$, which holds iff $\pi_\theta=q$.

\paragraph{Intuition.} MLE says “copy the data distribution.” Matching law says “allocate responses in proportion to reinforcement.” Both describe allocation matching contingencies.

\paragraph{Pitfalls.} Limited capacity/data, label smoothing, and regularization induce \emph{undermatching} (sensitivity $<1$) and \emph{bias}.

\subsection{Theorem 2: Generalized Matching via KL-Regularized Preferences}
\textbf{Claim.} Maximizing expected reward with a KL penalty to a reference policy yields a Boltzmann-rational form:
\begin{equation*}
\pi^*(y\given x)\propto \pi_{\mathrm{ref}}(y\given x)\,\exp\{\beta\, r(x,y)\}.
\end{equation*}
Hence, for any $y_1,y_2$,
\begin{equation*}
\log\frac{\pi^*(y_1\given x)}{\pi^*(y_2\given x)}
=\log\frac{\pi_{\mathrm{ref}}(y_1\given x)}{\pi_{\mathrm{ref}}(y_2\given x)}+\beta\,\big(r(x,y_1)-r(x,y_2)\big),
\end{equation*}
i.e., generalized matching with \emph{bias} from $\pi_{\mathrm{ref}}$ and \emph{sensitivity} $s=\beta$.

\paragraph{Objective.}
\begin{equation*}
\mathcal{J}(\pi)=\E_{y\sim\pi}[r(x,y)] - \beta\,\KL\!\big(\pi\;\Vert\;\pi_{\mathrm{ref}}\big).
\end{equation*}

\paragraph{Derivation (sketch).}
Lagrangian with normalization $\sum_y \pi(y)=1$:
\[
\mathcal{L}=\sum_y \pi(y)\,r(x,y) - \beta \sum_y \pi(y)\log\frac{\pi(y)}{\pi_{\mathrm{ref}}(y)} + \lambda\Big(\sum_y \pi(y)-1\Big).
\]
Set derivative w.r.t.\ $\pi(y)$ to zero and rearrange to obtain
$\pi(y)\propto \pi_{\mathrm{ref}}(y)\exp\{\beta r(x,y)\}$.
Taking log-odds between $y_1$ and $y_2$ gives the stated form.

\paragraph{Intuition.} The KL term keeps you near the reference (bias); reward scales your move by $\beta$ (sensitivity).

\paragraph{Pitfalls.} Reward misspecification, too-large $\beta$ (mode collapse), or highly peaked $\pi_{\mathrm{ref}}$ can dominate.

\subsection{Theorem 3: Temperature Scaling as Sensitivity Rescaling}
\textbf{Claim.} Sampling from $\pi_T(y\given x)\propto \pi(y\given x)^{1/T}$ yields
\begin{equation*}
\log\frac{\pi_T(y_1\given x)}{\pi_T(y_2\given x)}=\frac{1}{T}\log\frac{\pi(y_1\given x)}{\pi(y_2\given x)}.
\end{equation*}
Thus, temperature rescales \emph{sensitivity}: lower $T$ $\Rightarrow$ steeper allocation; higher $T$ $\Rightarrow$ flatter allocation.

\paragraph{Derivation.} Write $\pi_T(i)=\exp(z_i/T)/\sum_j \exp(z_j/T)$; subtract log-odds.

\paragraph{Pitfalls.} Very low $T$ may induce repetition; very high $T$ degrades coherence.

\subsection{Theorem 4: Mixtures and Regularization}
\textbf{Claim.} If $q=\sum_k \alpha_k\,q^{(k)}$ and the model has capacity, MLE yields $\pi^*=q$. Regularizers that shrink logits (label smoothing, entropy bonus, KL to $\pi_{\mathrm{ref}}$) produce \emph{undermatching} (slopes $<1$).

\paragraph{Sketch.} The $H+\KL$ decomposition still yields $\pi^*=q$. Odds under mixtures follow mixture odds; additive logit shrinkage reduces log-odds magnitudes (slopes $<1$).

\section{Key Takeaways from the Conclusions}
\begin{itemize}[leftmargin=*]
  \item MLE $\Rightarrow$ strict matching.  
  \item KL-regularized alignment $\Rightarrow$ generalized matching.  
  \item Temperature scaling $\Rightarrow$ sensitivity rescaling.  
  \item Allocation-based grounding: LLMs adapt like operant organisms, not like parrots.  
\end{itemize}

\section{Expanded Glossary with Context and Mathematics}

\paragraph{KL — Kullback–Leibler Divergence.}
\[
\KL(p\Vert q)=\sum_x p(x)\log\frac{p(x)}{q(x)}.
\]
\emph{Context:} In pretraining, minimizing cross-entropy $\Leftrightarrow$ minimizing $\KL(q\Vert \pi)$; in alignment, a $-\beta\,\KL(\pi\Vert \pi_{\mathrm{ref}})$ term discourages straying from a reference policy.

\paragraph{MLE — Maximum Likelihood Estimation.}
Choose parameters that maximize the probability of observed data; in our setting, equivalent to strict matching when realizable.

\paragraph{RLHF — Reinforcement Learning from Human Feedback.}
Learn from human preferences with reward + KL control; matches the generalized matching formulation.

\paragraph{DPO — Direct Preference Optimization.}
Align directly on preference pairs (no explicit reward model); still exhibits generalized matching dynamics via implicit KL control.

\paragraph{Cross-Entropy.}
\[
H(q,\pi)=-\sum_x q(x)\log \pi(x)=H(q)+\KL(q\Vert \pi).
\]
\emph{Context:} Core pretraining loss driving $\pi\to q$ (in the realizable limit).

\paragraph{Policy.}
A conditional distribution $\pi(y\given x)$ over outputs given inputs; LLMs can be analyzed as policies over tokens.

\paragraph{Logits.}
Unnormalized scores $z_i$ with softmax
\[
\pi(y=i\given x)=\frac{\exp(z_i)}{\sum_j \exp(z_j)}.
\]
\emph{Derivative:}
\[
\frac{\partial \pi(y=i)}{\partial z_k}=\pi(y=i)\big(\delta_{ik}-\pi(y=k)\big).
\]
\emph{CDF (for sampling):} $F(k)=\sum_{i\le k}\pi(y=i)$, and \emph{inverse CDF} sampling picks the smallest $k$ with $F(k)\ge u$, $u\sim U[0,1]$.\\
\emph{Context:} Logit differences control odds ratios, the natural unit for matching-law analysis.

\paragraph{Temperature.}
\[
\pi_T(i)=\frac{\exp(z_i/T)}{\sum_j \exp(z_j/T)}.
\]
\emph{Context:} Directly rescales sensitivity ($\propto 1/T$).

\paragraph{Sensitivity and Bias.}
From generalized matching (Baum): slope (\emph{sensitivity}) and intercept (\emph{bias}). In LLMs, sensitivity corresponds to $\beta$ (KL weight) or $1/T$; bias arises from $\pi_{\mathrm{ref}}$ or skewed data.

\paragraph{Grounding Types}
\begin{itemize}
\item \textbf{Allocation-grounding}
Grounding as lawful allocation to contingencies, not symbolic reference.  

\item \textbf{Gestalt-grounding} A form of grounding in which allocations are shaped by the integrated experiential substrate (\emph{gestalt}) available to the learner; for organisms this is their perceptual gestalt, while for LLMs it is the corpus as a distributed human gestalt, even in the absence of direct perception.
\end{itemize}

\paragraph{Operant semantics.}
Meaning emerges from stable allocations shaped by reinforcement.  
\paragraph{$n$-gram generator.}
Finite-order Markov chain text generator (e.g., \emph{Travesty}); parrots substrings, lacks global allocation laws.

\section{Broader Research Program (Expanded from Table~1)}
\label{sec:program}

This section lays out a practical research agenda that extends the paper’s formal correspondences into empirical and comparative work. It is organized around the rows in Table~\ref{tab:paradigm_mapping} and specifies core questions, minimal tests, and stretch goals. Cross–cutting reporting standards are included to enable cumulative comparisons across labs and model scales.

\subsection*{Cross–Cutting Metrics and Reporting}
For any experiment in this program, report: (i) \emph{sensitivity} (slope) and \emph{bias} (intercept) from log–odds regressions; (ii) $R^2$ with 95\% CIs; (iii) sample size and token/context granularity; (iv) any regularization (label smoothing, entropy bonuses, KL weights) and decoding controls (temperature, nucleus); (v) effective context length and any retrieval/external memory used.

\begin{table}[h!]
\centering
\begin{tabular}{|p{0.22\textwidth}|p{0.34\textwidth}|p{0.34\textwidth}|}
\hline
\textbf{Row (from Table~\ref{tab:paradigm_mapping})} & \textbf{Research Questions} & \textbf{Minimal Tests \& Stretch Goals} \\
\hline
Operant $\leftrightarrow$ MLE (strict matching) &
Over which context granularities (token, bigram, span) does strict matching hold? What are finite–sample or capacity–driven departures? &
\emph{Minimal:} Log–odds regressions of $\pi$ vs.\ corpus $q$ for token pairs across contexts; slope $\approx 1$, intercept $\approx 0$. \newline
\emph{Stretch:} Multi–token alternatives, compositional prompts, multilingual corpora; stratify by POS or semantic class.\\
\hline
Generalized matching $\leftrightarrow$ KL–regularized alignment &
How does KL weight $\beta$ map to sensitivity $s$ across datasets and model sizes? How large is reference–model bias? &
\emph{Minimal:} DPO with $\beta$ sweep on a public preference set; fit slope/bias from pre/post odds. \newline
\emph{Stretch:} PPO–style RLHF, multi–reference priors, human subject comparisons to animal sensitivity estimates.\\
\hline
Inference temperature $\leftrightarrow$ sensitivity rescaling &
Is the $1/T$ slope law stable across prompts, tasks, and truncation strategies (top–$k$, nucleus)? &
\emph{Minimal:} Temperature sweep; recover slopes $\approx 1/T$. \newline
\emph{Stretch:} Adaptive temperature controllers; task–conditioned sensitivity and uncertainty calibration.\\
\hline
Classical conditioning $\leftrightarrow$ supervised learning &
Do cue–outcome learning curves mirror associative models (e.g., blocking, overshadowing)? &
\emph{Minimal:} Train small nets on synthetic cue–outcome pairs; fit Rescorla–Wagner–like curves. \newline
\emph{Stretch:} Pretraining–slice analyses for multi–cue effects; ablations isolating “blocking–like” behavior.\\
\hline
Hebbian plasticity $\leftrightarrow$ self–supervised representation learning &
Are embedding updates consistent with correlation–driven rules? &
\emph{Minimal:} Track embedding drift under masked/causal LM objectives; correlate update magnitude with co–occurrence. \newline
\emph{Stretch:} Compare with neurophys data (STDP–like signatures), probe layerwise locality vs.\ global updates.\\
\hline
Observational learning $\leftrightarrow$ in–context learning &
How do demonstrations in context shape one–shot generalization and allocation fidelity? &
\emph{Minimal:} Vary number/type of in–context examples; measure performance and odds–slope changes. \newline
\emph{Stretch:} Derive sample–efficiency laws; compare patterns to animal imitation/observational datasets.\\
\hline
Memory constraints $\leftrightarrow$ context windows &
How does allocation fidelity scale with context length and recency? What is the effect of external memory? &
\emph{Minimal:} Strict–matching regressions across window sizes; quantify undermatching vs.\ distance. \newline
\emph{Stretch:} Add retrieval–augmented generation as “extended working memory”; estimate gains as sensitivity restoration.\\
\hline
\end{tabular}
\caption{Expanded research program connecting behavior–analytic constructs to LLM mechanisms. Each row specifies core questions, a minimal test, and stretch goals, with shared reporting standards for sensitivity, bias, $R^2$, and capacity/recency controls.}
\label{tab:program}
\end{table}

\subsection*{Core Tracks and Milestones}
\paragraph{Track A: Allocation Benchmarks (3–6 weeks).}
Replicate strict matching (MLE), generalized matching (KL–regularized alignment), and temperature scaling on small models (e.g., 124M–355M). Milestone: slopes near target values with clear CIs; bias quantified.

\paragraph{Track B: Capacity and Memory (4–8 weeks).}
Repeat Track A while varying context window and model size. Milestone: sensitivity improves with window/scale; undermatching increases with distance.

\paragraph{Track C: Comparative Phenomena (6–10 weeks).}
Run supervised cue–outcome and in–context learning studies; test associative predictions and demonstration effects. Milestone: recover classical conditioning curves; quantify observational–learning analogs.

\paragraph{Track D: Mechanistic Probes (optional).}
Analyze embedding dynamics and layerwise contributions to allocation changes during alignment/fine–tuning.

\subsection*{Risks and Mitigations}
\emph{Data leakage/bias:} stratify by domain and control for frequency. \newline
\emph{Overfitting to proxies:} validate on held–out corpora and alternate tokenizations. \newline
\emph{Confounded temperature/filters:} report exact decoding stack (top–$k$, $p$, repetition penalties).

\subsection*{Artifacts and Reproducibility}
Release code, seeds, and config files; provide CSVs with per–pair odds, slopes, intercepts, and CIs; include prompts and exact model hashes/versions to facilitate replication.

\section*{Big Picture: From Pigeons to Programs}

The matching law was first discovered in experiments with pigeons in the 1960s
and refined by later behavioral scientists. Far from being outdated, it is still
used today in applied behavior analysis, neuroscience, and even everyday
contexts like how people divide attention when texting. In this guide, we have
shown how the same principles can also be found in large language models. 

By thinking in terms of both \emph{allocation-grounding} (the general law that
allocations track contingencies) and \emph{gestalt-grounding} (the specific
channel through which contingencies are experienced), we can see a continuous
thread: from animals in Skinner boxes, to human decision-making, to today’s
most advanced AI models. What once seemed like separate domains of learning are
connected by the same underlying laws.


\section{Micro-Exercises (for Intuition)}
\begin{enumerate}[leftmargin=*]
  \item \textbf{Strict Matching.} If $q(a)=0.8,q(b)=0.2$, then $\pi^*(a)=0.8,\pi^*(b)=0.2$.  
  \item \textbf{Generalized Matching.} With $\pi_{\mathrm{ref}}$ uniform, $r(y_1)=2$, $r(y_2)=1$, $\beta=1$, odds ratio = $e^1\approx 2.7$.  
  \item \textbf{Temperature.} If log-odds = 2 at $T=1$, then =1 at $T=2$.  
  \item \textbf{Mixture.} If half from $q^{(1)}(a)=0.9$, half from $q^{(2)}(a)=0.1$, then $q(a)=0.5$.  
  \item \textbf{$n$-gram vs.\ LLM.} Run a toy $n$-gram generator and a small LLM on the same seed. Compare outputs: where does parroting end and adaptive allocation begin?  
\end{enumerate}

\section*{Build Notes}
This document uses \texttt{amsmath}. If you copy equations elsewhere, keep text like “(strict matching)” \emph{outside} math, or wrap it with \verb|\text{}| (requires \texttt{amsmath}). Avoid literal \verb|{[}| \verb|{]}|; use parentheses or \verb|\big[ \big]| in math.

\end{document}
